% ---------------------------------------------------------------------------
% Data window memo for supervisor consultation.
% Compiles with pdflatex (2 passes for the table of contents / references).
%
% LATEX NOTES, learned the hard way on the earlier progress-report conversion:
%
%   1. \textcent is a TEXT-mode command. Used bare inside $...$ it emits only
%      a WARNING (not an error) and silently drops the glyph, so "$4.3\cent$"
%      renders as "4.3". Wrapping it in \text{} (amsmath) makes it valid in
%      both modes. Never use \textcent directly in maths.
%   2. Maths inside \textbf{} is NOT bolded. Use \bm{} from the bm package for
%      bold symbols and numbers inside maths.
%   3. When checking the build, count WARNINGS as well as errors. A clean
%      "0 errors" hid a dropped-glyph bug once already.
% ---------------------------------------------------------------------------
\documentclass[11pt,a4paper]{article}

\usepackage[T1]{fontenc}
\usepackage[utf8]{inputenc}
\usepackage{textcomp}
\usepackage{amsmath}
\usepackage{bm}
\usepackage{booktabs}
\usepackage{longtable}
\usepackage{siunitx}
\usepackage[table]{xcolor}
\usepackage[margin=2.4cm]{geometry}
\usepackage{enumitem}
\usepackage{url}
\usepackage{hyperref}
\hypersetup{colorlinks=true, linkcolor=black, urlcolor=blue!60!black}

% Cent sign valid in BOTH text and math mode. See note 1 above.
\newcommand{\cent}{\text{\textcent}}

\setlist[itemize]{itemsep=2pt, topsep=4pt}
\setlist[enumerate]{itemsep=3pt, topsep=4pt}
\sisetup{group-separator={,}, group-minimum-digits=4, detect-weight=true}

\title{\vspace{-1.6cm}\textbf{Data window for the MECE STGAT experiment}\\[2pt]
\large What 54 months of Kalshi data actually contains}
\author{}
\date{\vspace{-0.9cm}26 September 2026}

\begin{document}
\maketitle
\vspace{-1.2cm}

\section*{Summary and the ask}

You raised that twelve months is not enough data. The archive does span
2021-06 to 2025-11, which is 54 months, but the usable portion is much
smaller and I have now measured it rather than assumed it. This memo sets out
the measurement, the three constraints it reveals, and the split design I
propose instead of a longer calendar. Four questions where your view would
change what I do are at the end (Section~\ref{sec:questions}).

\begin{itemize}
  \item Ten of the 54 months contain \emph{zero} trades on any candidate
        basket leg. This is a hard floor.
  \item A 23-month middle period averages 37 usable basket-days per month at
        a 5.5-hour median gap between legs, which is stale pricing rather
        than mispricing.
  \item Platform traded notional grew by a factor of 369 over the archive, so
        pooling across it estimates a parameter that exists in neither
        regime.
  \item I propose 13 months (2024-11 to 2025-11) with seven expanding-window
        walk-forward folds, giving seven forward test evaluations instead of
        one.
  \item Separately, and since we last spoke: a measurement check has shown
        that 92.5\% of the apparent basket deviation is the bid/ask spread
        rather than mispricing, so I have withdrawn the profitability claim.
        Section~\ref{sec:lit} sets out the test. The forecasting results are
        unaffected. This lowers what the data window has to support, and I
        would value your view on the reframe (question~6).
\end{itemize}

\section{Why months are the wrong unit}

The mechanism under test is the sum-to-one constraint on a mutually exclusive,
collectively exhaustive basket: if exactly one leg of a partition resolves
\textsc{yes}, the legs' prices must sum to \$1 at resolution. Measuring a
deviation therefore requires a price for \emph{every} leg. If two of six legs
have not traded, there is no basket sum, no deviation, and nothing to trade or
forecast, however long the calendar is.

The unit of available data is therefore the \textbf{fully covered
basket-day}: one (event, day) observation where every declared leg traded. A
month with 12 such observations is not a month of data.

\section{Measured availability}

Identification follows the resolved-outcome test already used in the project
(fully resolved, exactly one \textsc{yes} leg, not a threshold ladder, not a
combination proposition, and family-consistent across at least three
historical instances). Coverage is then counted per month. Of 54 months, 44
have at least one fully covered basket, 41 have at least ten, and 21 have at
least one hundred. Grouped:

\begin{table}[ht]
\centering\small
\begin{tabular}{lrrrrr}
\toprule
Period & Months & Basket-days & Per mth & Median $|\mathrm{dev}|$ & Median gap \\
\midrule
2021-06 to 2022-03 & 10 & 0 & 0 & --- & --- \\
2022-04 to 2022-11 & 8 & 914 & 114 & $3.1\cent$ & 3.15\,h \\
2022-12 to 2024-10 & 23 & 861 & 37 & $8.1\cent$ & 5.49\,h \\
\textbf{2024-11 to 2025-11} & \textbf{13} & \textbf{5,806} & \textbf{447} & $\bm{4.3\cent}$ & \textbf{3.95\,h} \\
\midrule
\emph{current pilot} 2025-05 to 09 & \emph{5} & \emph{3,242} & \emph{648} & $3.8\cent$ & 3.09\,h \\
\bottomrule
\end{tabular}
\caption{Fully covered MECE basket-days per period. ``Median leg gap'' is the
span between the earliest and latest last-trade time among a basket's legs on
the day it was counted as covered.}
\end{table}

\section{Three constraints, each measured}

\subsection{The first ten months have no coverage at all}

2021-06 to 2022-03 contain no trades on any leg of any candidate basket. Not
thin coverage: none. This is a property of Kalshi's early catalogue, not a
methodological choice, and no reanalysis recovers those months.

\subsection{The middle 23 months are stale, not mispriced}
\label{sec:stale}

The 2022-12 to 2024-10 period averages 37 basket-days per month, against 447
in the proposed window. More importantly the median gap between the first and
last leg trade on a ``covered'' day rises to 5.49 hours, and the apparent
deviation rises with it:

\begin{table}[ht]
\centering
\begin{tabular}{lrr}
\toprule
Period & Median $|\mathrm{dev}|$ & Median leg gap \\
\midrule
2022-04 to 2022-11 & $3.1\cent$ & 3.15\,h \\
2023-07 to 2023-12 & $5.9\cent$ & 5.4\,h \\
2024-01 to 2024-06 & $13.2\cent$ & 5.6\,h \\
2024-11 to 2025-11 & $4.3\cent$ & 3.95\,h \\
\bottomrule
\end{tabular}
\caption{Apparent deviation scales with the leg gap, not with market
inefficiency. The 2024-01 to 2024-06 rows peak at $15.5\cent$ at gaps of 5 to
6.6 hours.}
\end{table}

When a basket's legs last traded a median of 5.6 hours apart, most of a
$13\cent$ ``deviation'' is the market having moved between the legs. Training
on that period would teach staleness as signal and would inflate the
reported opportunity count with artefacts. I propose to exclude it from
fitting and to cite it only as a data-availability figure.

\subsection{A 369-fold liquidity regime change}

Monthly traded notional across the archive:

\begin{table}[ht]
\centering
\begin{tabular}{lrr}
\toprule
Month & Total notional & Weather/climate notional \\
\midrule
2024-06 & \$3.2M & \$0.31M \\
2024-09 & \$3.7M & \$0.54M \\
2024-12 & \$24.4M & \$10.3M \\
2025-05 & \$211.0M & \$2.95M \\
2025-09 & \$490.9M & \$4.42M \\
2025-11 & \$1{,}193M & \$4.77M \\
\bottomrule
\end{tabular}
\caption{From 2024-06 to 2025-11 total notional grows by a factor of 369. The
current five-month pilot averages \$264M per month; 2024-06 to
2024-09 averages \$3.75M, a factor of 70.}
\end{table}

Every quantity the experiment measures, namely the implied per-leg spread,
basket completeness, deviation persistence and executable participation size,
is a function of market depth. Pooling observations across a 70-fold depth
change estimates a parameter that does not exist in either regime. The window
is therefore bounded by \textbf{liquidity stationarity} rather than by
calendar availability. Within the proposed 13 months the growth is still a
factor of roughly 24, which I will report as a stated limitation.

\section{A composition break as well as a thinning}

Sports accounts for 51.9\% of the current basket population (24.4\% of
opportunities above $5\cent$) and did not exist before 2025-01. Kalshi could
not list sporting event contracts until the regulatory position shifted in
late 2024; the first sports markets appear in 2024-12 and scale through Q1
2025. An independent ticker-pattern sweep that does not consult the category
field confirms only eight sports-pattern markets before 2024-12, and all eight
are news-about-sports contracts (NFL viewership, the MLB collective
bargaining agreement, the Oakland relocation) rather than wagers on outcomes.

Consequently the 2022 baskets are 81\% to 94\% weather. Adding them would not
merely add thinner data, it would mix a different population into the same
estimate.

\section{Proposed design}

\subsection{Primary: expanding-window walk-forward, strictly forward in time}

Use 2024-11 to 2025-11 (13 months, one liquidity regime, 1.8 times the
basket-days of the current pilot) and replace the single fixed split with
seven expanding-window folds:

\begin{table}[ht]
\centering
\begin{tabular}{clll}
\toprule
Fold & Train & Validate & Test \\
\midrule
1 & 2024-11 to 2025-03 & 2025-04 & 2025-05 \\
2 & 2024-11 to 2025-04 & 2025-05 & 2025-06 \\
3 & 2024-11 to 2025-05 & 2025-06 & 2025-07 \\
4 & 2024-11 to 2025-06 & 2025-07 & 2025-08 \\
5 & 2024-11 to 2025-07 & 2025-08 & 2025-09 \\
6 & 2024-11 to 2025-08 & 2025-09 & 2025-10 \\
7 & 2024-11 to 2025-09 & 2025-10 & 2025-11 \\
\bottomrule
\end{tabular}
\caption{Every fold trains only on its own past, so the evaluation remains a
forecast throughout.}
\end{table}

This yields seven forward test evaluations rather than one, so the PnL and the
model-versus-control gap become distributions with reportable dispersion
instead of single point estimates. It also covers a full winter (2024-11 to
2025-03), which the summer-only pilot could not; since weather is 58\% of the
opportunity population above $5\cent$, that is what separates ``the model does
not generalise'' from ``a different season is a different market''.

\subsection{Secondary: 2022 as a clearly labelled backcast}

The 2022-04 to 2022-11 block retains value as a structural transfer check,
reported separately and never merged into the headline metric:

\begin{quote}
The model trained on 2024-11 to 2025-11 was additionally applied to 2022-04 to
2022-11, a period 26 months before the training window, on a venue roughly 60
times thinner and under the previous ticker convention. This is a
\emph{backcast} and is not evidence of forecasting performance; it tests only
whether the learned basket structure holds in an earlier liquidity regime.
\end{quote}

I want to be explicit that using 2022 as a \emph{test set} would be wrong,
because the training data would postdate it. That is why it appears here as a
labelled backcast rather than as the held-out split.

\section{Effective sample size: two different questions}

\begin{itemize}
  \item \textbf{For fitting, the data is not scarce.} The model trains on
        hourly graph snapshots across 500 to 1,100 active nodes per timestep,
        roughly 9,000 snapshots in the proposed window. The forecast head is
        zero-initialised, so the model begins at persistence rather than at
        noise.
  \item \textbf{For inference, the effective unit is the basket-day and it is
        modest.} Confidence intervals are cluster-bootstrapped by basket, so
        the effective $n$ is the number of distinct baskets (5,806 in the
        proposed window), not the roughly 46,000 basket-hours. Consecutive
        hourly snapshots of one basket are largely frozen and are not
        independent observations.
\end{itemize}

On current evidence the primary comparison is adequately powered on the
five-month pilot alone: the model beats its randomised-timing control by
\$0.170 per trade with a cluster-bootstrapped 95\% interval of
$[+0.107, +0.247]$. Extending to 13 months and seven folds is therefore not
required for statistical significance. It is required for
\emph{generalisation}, which I take to be the substance of your concern.

\section{Comparison with the published literature}
\label{sec:lit}

Since the question is how much data this phenomenon requires, I checked the
sample periods of the three papers nearest to it. An important qualification
first: \textbf{only one of the three treats the sum-to-one constraint as a
primary object of study}. In the other two it appears as a minor category
alongside a different mechanism, and in both it turns out to be nearly empty.
None of the three trains a forecasting model, so none has a train/test split.

\begin{table}[ht]
\centering\small
\begin{tabular}{lllll}
\toprule
Study & Window & Length & Sum-to-one component & Model \\
\midrule
Unravelling the Probabilistic & 1 Apr 2024 to & 12 months & \textbf{co-equal primary} & none \\
\quad Forest (AFT 2025) & 1 Apr 2025 & & type (``market & \\
 & & & rebalancing'') & \\[3pt]
Arbitrage Analysis in Poly- & 4 Feb 2026 to & 1 month & marginal: 7 episodes & none \\
\quad market NBA Markets & 4 Mar 2026 & & across 3,042 markets, & \\
 & & & vs 290 combinatorial & \\[3pt]
Executable Arbitrage and & 14 Apr 2026 to & 5 weeks & marginal: $\sim$3\% of & none \\
\quad Market Efficiency & 19 May 2026 & & value; 97\% from a & \\
 & & & converter primitive & \\
\midrule
\textbf{This project (proposed)} & \textbf{2024-11 to} & \textbf{13 months} & \textbf{sole mechanism} & STGAT, \\
 & \textbf{2025-11} & & & 7 folds \\
\bottomrule
\end{tabular}
\caption{The three nearest studies, all on Polymarket. The third study's
window is its level-2 order book panel, hour-stratified; its transaction
panel covers events listed through 31 December 2025.}
\end{table}

Three observations:

\begin{enumerate}
  \item \textbf{The only study that treats sum-to-one as a primary type uses
        12 months}, which the proposed window slightly exceeds. Its other
        type, combinatorial arbitrage, spans semantically dependent markets
        and is a different mechanism from this project's.
  \item \textbf{The two short-window studies are not evidence about how long
        a sum-to-one study needs to be.} Their sum-to-one components are
        incidental and nearly empty: 7 episodes across 3,042 markets in one,
        roughly 3\% of measured value in the other. Their windows were chosen
        for the mechanisms that dominate their results, not for this one.
  \item \textbf{None of the three trains a model}, so none faces the
        train/test ordering constraint. The NBA paper states explicitly that
        ``there is no train/test split because the analysis doesn't involve
        machine learning or out-of-sample validation''.
\end{enumerate}

I want to be careful not to over-claim. A detection study needs enough
observations to characterise a distribution; a forecasting study needs enough
to fit a model \emph{and} evaluate it out of sample, which is strictly
harder. So ``12 months is the published maximum'' does not establish that 13
months suffices for training. What it establishes is that there is no
precedent in this literature for requiring multi-year windows, and that the
one genuinely comparable study used slightly less than the proposed window.

\paragraph{A comparison I am deliberately not making.} It is tempting to set
this project's 7,581 fully covered basket-days against the NBA paper's 7
arbitrage episodes and conclude the mechanism is more present on Kalshi.
That comparison is not valid, and I would rather say so than have it
challenged.

The two quantities are different objects. The NBA paper detects instants at
which posted \emph{asks} satisfied $\mathrm{Ask}_A + \mathrm{Ask}_B < \$1$,
polling every 3.6 to 5.5 seconds: a strictly executable condition at a single
moment. This project counts basket-days on which every leg traded and the
\emph{last trade prices} did not sum to \$1, with those prints separated by a
median of 2.3 to 6 hours. A last trade is struck at the bid or the ask, never
the mid, so for a six-leg basket with a per-leg spread near $5.6\cent$ the sum
of mixed last prints has a standard deviation of roughly
$\sqrt{6}\times2.8\cent \approx 6.9\cent$ around \$1 with no mispricing at all,
whose median absolute value is about $4.6\cent$. The measured median is
$4\cent$.

The bulk of the deviation population is therefore consistent with bid/ask
mixing rather than mispricing, and I have now tested that directly rather
than leaving it as a caveat.

Kalshi trades carry \texttt{taker\_side}: a taker buying \textsc{yes} lifts
the ask, one buying \textsc{no} sells \textsc{yes} into the bid. Across one
month, 68.7\% of prints are ask-side. With a half-spread of $2.8\cent$ that
predicts a signed mean deviation of $+6.3\cent$ on a six-leg basket from
order-flow imbalance alone. \textbf{The prediction was recorded before the
measurement. Measured: $\bm{+5.25\cent}$, with 92.5\% of 4{,}891 baskets
summing above \$1} --- the classic betting-market overround, where a genuine
mispricing population would be roughly balanced either side of \$1.

Pricing each leg on its own side, restricted to baskets whose legs all
printed within an hour of each other:

\begin{center}\small
\begin{tabular}{lr}
\toprule
buy every leg at the \textbf{ask} & \$1.1334 \\
sell every leg at the \textbf{bid} & \$0.9993 \\
a basket settles at & \$1.0000 \\
\bottomrule
\end{tabular}
\end{center}

\noindent Selling the apparently overpriced basket at executable prices nets
$-0.07\cent$ gross, before any fee; Kalshi's taker fee on a six-leg basket is
$1.3\cent$ to $5.8\cent$ depending on the leg prices. Restricting to the
large-deviation tail does not help: the share of baskets with a positive
executable edge \emph{falls} from 84\% to 54\%, because a large apparent
deviation is produced by a wide spread, and wide-spread baskets execute
worse.

\textbf{I am therefore withdrawing the profitability claim} and reporting
this as the finding in its own right: a backtest on prediction-market
partitions that prices legs from trade prints without regard to execution
side manufactures profit that does not exist at executable prices. The
forecasting results are unaffected, since they concern price accuracy rather
than realised trading, and the model-versus-control comparison is a paired
test on identical baskets that is provably invariant to the cost model.

It is also worth noting that the AFT window (April 2024 to April 2025) spans
Polymarket's election-driven liquidity expansion, a regime change comparable
in kind to the one in Section~3.3, and that the paper reports results pooled
across it rather than segmented by regime.

\section{Questions}
\label{sec:questions}

\begin{enumerate}
  \item Does seven-fold expanding walk-forward within a single liquidity
        regime address your concern, or is the concern specifically calendar
        span? These have different remedies and I would rather solve the one
        you mean.
  \item Would you accept 2022-04 to 2022-11 as a clearly labelled backcast
        for structural transfer, reported separately from the headline?
  \item Would you prefer I include the thin 2022-12 to 2024-10 period with an
        explicit staleness filter, accepting noisier labels, rather than
        excluding it? Section~\ref{sec:stale} is my argument against, but I
        would rather it were your decision than my omission.
  \item Is the basket-day the right unit to report as effective sample size,
        given that the hourly snapshots within a basket are largely frozen?

  \item Section~\ref{sec:lit} finds that the only comparable study of
        sum-to-one arbitrage uses 12 months, and that none of the three
        nearest studies trains a model or uses a train/test split. Is there a
        departmental expectation on window length that I should be matching
        instead of the literature's, and if so what is it? I would rather
        know the target than infer it.

  \item Given the withdrawn profitability claim, the project's contribution
        becomes (a) the measurement finding that last-trade deviations on
        these partitions are dominated by the overround, (b) forecasting
        accuracy against persistence and an optimal shrinkage oracle, with an
        ablation attributing roughly 92\% of the margin to the graph, and (c)
        the ladder mechanism as a diagnosed negative control. Is that an
        acceptable shape for the dissertation, or would you want the emphasis
        elsewhere? The data window matters less under this framing, since it
        now has to support forecasting folds rather than a trading sample.
\end{enumerate}

\section*{Reproduction}

All figures above come from two scripts, both with built-in self-tests:

\begin{itemize}
  \item \path{stg_infra/examples/mece_yield_by_month.py}, which pools
        basket identification across all 54 months (the family-consistency
        test needs three instances, which may span a year) and then counts
        full-basket coverage per trades month. Output is saved to
        \path{mece_yield_by_month.parquet}.
  \item \path{stg_infra/examples/sports_categorisation_check.py}, which
        verifies the sports timeline three ways: a ticker-pattern sweep
        independent of the category field, a check that the category lookup
        is not time-bounded, and a trading-volume breakdown.
\end{itemize}

The full per-month counts are in Appendix~\ref{app:permonth} and the sports
timeline checks in Appendix~\ref{app:sports}.

\appendix

\section{Per-month MECE basket coverage}
\label{app:permonth}

Every month in 2021-06 to 2022-03 is absent because it contained no trades on
any candidate basket leg. ``Partial'' counts (event, day) rows with at least
one leg trading but not all of them, so the ratio of the first two columns is
the coverage rate. ``Top category'' is the share of that month's fully covered
baskets in the single most common category.

{\footnotesize
\begin{longtable}{lrrrlrr}
\toprule
Month & Full & Partial & Fam. & Top category & Med. $|\mathrm{dev}|$ & Med. gap \\
\midrule
\endfirsthead
\toprule
Month & Full & Partial & Fam. & Top category & Med. $|\mathrm{dev}|$ & Med. gap \\
\midrule
\endhead
\midrule \multicolumn{7}{r}{\emph{continued on next page}} \\
\endfoot
\bottomrule
\endlastfoot
2022-04 & 21 & 26 & 5 & weather 81\% & $3.0\cent$ & 3.11 \\
2022-05 & 137 & 227 & 5 & weather 81\% & $3.0\cent$ & 3.03 \\
2022-06 & 120 & 220 & 4 & weather 93\% & $2.0\cent$ & 2.75 \\
2022-07 & 119 & 247 & 4 & weather 93\% & $3.0\cent$ & 2.99 \\
2022-08 & 125 & 259 & 3 & weather 94\% & $3.0\cent$ & 3.10 \\
2022-09 & 133 & 264 & 4 & weather 83\% & $3.0\cent$ & 3.47 \\
2022-10 & 129 & 257 & 5 & weather 88\% & $4.0\cent$ & 3.12 \\
2022-11 & 130 & 247 & 3 & weather 92\% & $4.0\cent$ & 3.60 \\
\midrule
2022-12 & 60 & 261 & 3 & weather 92\% & $3.0\cent$ & 4.88 \\
2023-01 & 58 & 259 & 3 & weather 76\% & $5.0\cent$ & 5.93 \\
2023-02 & 37 & 223 & 3 & weather 76\% & $5.0\cent$ & 6.06 \\
2023-03 & 39 & 251 & 4 & weather 72\% & $5.0\cent$ & 6.22 \\
2023-04 & 49 & 248 & 4 & weather 67\% & $8.0\cent$ & 6.70 \\
2023-05 & 28 & 303 & 4 & weather 71\% & $7.5\cent$ & 5.41 \\
2023-06 & 54 & 367 & 6 & weather 52\% & $6.0\cent$ & 5.97 \\
2023-07 & 18 & 367 & 6 & weather 61\% & $13.5\cent$ & 5.99 \\
2023-08 & 15 & 390 & 3 & weather 100\% & $10.0\cent$ & 4.82 \\
2023-09 & 7 & 341 & 2 & weather 100\% & $2.0\cent$ & 5.92 \\
2023-10 & 19 & 333 & 4 & weather 95\% & $5.0\cent$ & 2.36 \\
2023-11 & 8 & 332 & 3 & weather 62\% & $4.5\cent$ & 7.05 \\
2023-12 & 8 & 329 & 3 & weather 62\% & $5.5\cent$ & 6.29 \\
2024-01 & 17 & 318 & 3 & weather 88\% & $11.0\cent$ & 2.78 \\
2024-02 & 16 & 347 & 4 & weather 94\% & $13.5\cent$ & 6.20 \\
2024-03 & 34 & 384 & 5 & weather 76\% & $15.5\cent$ & 5.21 \\
2024-04 & 58 & 391 & 5 & weather 84\% & $11.5\cent$ & 5.15 \\
2024-05 & 49 & 406 & 6 & weather 96\% & $14.0\cent$ & 6.63 \\
2024-06 & 58 & 441 & 7 & weather 83\% & $13.5\cent$ & 6.43 \\
2024-07 & 100 & 492 & 7 & weather 91\% & $7.0\cent$ & 5.86 \\
2024-08 & 64 & 450 & 5 & weather 97\% & $8.0\cent$ & 5.67 \\
2024-09 & 28 & 423 & 4 & weather 100\% & $6.5\cent$ & 6.13 \\
2024-10 & 37 & 454 & 8 & weather 100\% & $5.0\cent$ & 2.62 \\
\midrule
\textbf{2024-11} & \textbf{332} & 704 & 9 & weather 97\% & $4.5\cent$ & 4.20 \\
\textbf{2024-12} & \textbf{274} & 1{,}139 & 8 & weather 97\% & $4.0\cent$ & 4.21 \\
\textbf{2025-01} & \textbf{241} & 1{,}354 & 11 & weather 94\% & $5.0\cent$ & 5.08 \\
\textbf{2025-02} & \textbf{312} & 1{,}207 & 12 & weather 91\% & $6.0\cent$ & 4.99 \\
\textbf{2025-03} & \textbf{346} & 1{,}647 & 14 & weather 73\% & $5.0\cent$ & 5.87 \\
\textbf{2025-04} & \textbf{256} & 1{,}293 & 12 & weather 82\% & $5.0\cent$ & 5.13 \\
\textbf{2025-05} & \textbf{331} & 1{,}491 & 16 & weather 66\% & $5.0\cent$ & 4.29 \\
\textbf{2025-06} & \textbf{590} & 1{,}710 & 13 & sports 49\% & $4.0\cent$ & 2.80 \\
\textbf{2025-07} & \textbf{649} & 1{,}979 & 15 & sports 53\% & $4.0\cent$ & 3.44 \\
\textbf{2025-08} & \textbf{803} & 2{,}522 & 16 & sports 63\% & $3.0\cent$ & 2.63 \\
\textbf{2025-09} & \textbf{869} & 2{,}565 & 15 & sports 73\% & $3.0\cent$ & 2.28 \\
\textbf{2025-10} & \textbf{530} & 1{,}596 & 17 & weather 58\% & $4.0\cent$ & 2.65 \\
\textbf{2025-11} & \textbf{273} & 870 & 10 & weather 85\% & $4.0\cent$ & 3.81 \\
\midrule
\multicolumn{2}{l}{\textbf{Total, 44 months}} & \multicolumn{5}{l}{7,581 fully covered basket-days} \\
\end{longtable}}

\noindent Bold rows are the proposed window. The horizontal rules separate the
three periods of Table~1. Note how the median leg gap tracks the coverage
count: the thin 2023--2024 rows sit at 5 to 7 hours while both dense periods
sit near 3 to 4, which is the evidence for Section~\ref{sec:stale}.

\section{Sports timeline checks}
\label{app:sports}

The claim that sports did not exist before 2025-01 was checked three ways,
because a categorisation artefact would have undermined the window argument.
The published figure that sports is over 90\% of Kalshi activity in 2025
initially appeared to contradict a 2024-12 first-seen month.

\subsection*{B.1 Ticker-pattern sweep, independent of the category field}

Series prefixes were matched against sports patterns without consulting
Kalshi's category field or the category lookup file. Patterns were split into
\emph{strong} (unambiguous: NFL, NBA, NASCAR, LALIGA, \ldots) and \emph{weak}
(ambiguous: GAME, WINS, CHAMPION, \ldots), because election series such as
\path{KXTRUMPWINS} match the latter and would otherwise fabricate a
pre-2024-12 sports population.

\begin{table}[ht]
\centering\small
\begin{tabular}{lrr@{\hskip 1.2cm}lrr}
\toprule
Month & Strong & Weak & Month & Strong & Weak \\
\midrule
2021-09 & 1 & --- & 2025-03 & 1,362 & 143 \\
2021-10 & 3 & 1 & 2025-04 & 1,154 & 312 \\
2021-11 & 2 & --- & 2025-05 & 2,969 & 170 \\
2023-06 & 1 & --- & 2025-06 & 2,896 & 842 \\
2023-07 & 1 & --- & 2025-07 & 3,430 & 1,125 \\
2024-04 & --- & 11 & 2025-08 & 8,770 & 1,210 \\
2024-06 & --- & 1 & 2025-09 & 11,425 & 17,574 \\
2024-11 & --- & 50 & 2025-10 & 9,165 & 120,311 \\
2024-12 & 1 & 5 & 2025-11 & 7,023 & 81,791 \\
2025-01 & 116 & --- & & & \\
2025-02 & 585 & 160 & & & \\
\bottomrule
\end{tabular}
\caption{Sports-pattern market counts by month. Months absent from the table
matched nothing. The weak column explodes from 2025-09 with the
\texttt{exotics} NFL and NBA parlay products.}
\end{table}

Eight markets across three series match before 2024-12, and none is a wager on
a sporting outcome:

\begin{table}[ht]
\centering\small
\begin{tabular}{llrl}
\toprule
Series & First month & Markets & Kalshi category \\
\midrule
\texttt{NFLVIEW} & 2021-09 & 4 & entertainment \\
\texttt{MLBOAK} & 2021-11 & 3 & entertainment \\
\texttt{MLBCBA} & 2021-11 & 1 & entertainment \\
\bottomrule
\end{tabular}
\caption{NFL viewership, the Oakland relocation and the MLB collective
bargaining agreement: news about sports, not outcome contracts.}
\end{table}

\subsection*{B.2 The category lookup is not time-bounded}

The suspected artefact was that \path{kalshi_series_categories.parquet} might
map only currently listed series, sending retired ones to a nine-word keyword
fallback and under-categorising early months by construction. It does not: the
lookup maps 13,738 series, and the share of each month's series that it
resolves is $1.00$ for every month of 2021, with a minimum of $0.86$
(2022-09) across all 54 months. Only 16 series in the entire archive are
unmapped, of which one is sports (\texttt{KXMLBWINS}, 150 markets, first seen
2025-05).

\subsection*{B.3 Share of trading activity, not of listings}

The apparent contradiction dissolves once the two quantities are separated.
The audit that produced the 2024-12 first-seen month counts \emph{distinct
markets listed}; the published figure is a claim about \emph{traded volume}.
Both hold simultaneously, because crypto and financials mint hundreds of
thousands of hourly strike markets that barely trade, while one game market
absorbs enormous size. Sports is roughly 2\% of listed markets (60,243 of
about 4.2 million) and the following share of traded notional:

\begin{table}[ht]
\centering\small
\begin{tabular}{lrr@{\hskip 1.2cm}lrr}
\toprule
Month & Total notional & Sports & Month & Total notional & Sports \\
\midrule
2024-06 & \$3.2M & --- & 2025-04 & \$136.9M & 69.8\% \\
2024-09 & \$3.7M & --- & 2025-05 & \$211.0M & 78.5\% \\
2024-11 & \$49.6M & --- & 2025-06 & \$230.2M & 83.4\% \\
2024-12 & \$24.4M & --- & 2025-07 & \$163.4M & 75.2\% \\
2025-01 & \$31.1M & 4.1\% & 2025-08 & \$226.4M & 80.0\% \\
2025-02 & \$26.0M & 3.8\% & 2025-09 & \$490.9M & 82.3\% \\
2025-03 & \$187.3M & \textbf{86.9\%} & 2025-10 & \$1{,}073.0M & 86.5\% \\
 & & & 2025-11 & \$1{,}193.0M & \textbf{87.2\%} \\
\bottomrule
\end{tabular}
\caption{Sports is absent from the trades files before 2025-01, then takes
over in 2025-03. The 87.2\% in 2025-11 reproduces the published revenue-share
figure from this project's own data.}
\end{table}

\paragraph{A limitation of the sweep.} In 2025-03 the category field puts
sports at 86.9\% of notional while the keyword sweep of B.1 sees 0.7\%. The
sweep under-detects sports by two orders of magnitude in that month, which is
the characteristic failure of keyword heuristics and the reason B.2 matters:
the audit never depended on keywords. B.1 should be read as corroboration, not
as the primary evidence.

\end{document}
